\chapter{Discusiones, conclusiones y futuro}\label{Capitulo_4}
\chapterquote{}{}

\section{Discusiones}\label{IV:discusiones}

El objetivo de esta tesis fue reproducir con un modelo de caminante estocástico el movimiento diario de las tortugas registrado con GPS, y describir
cómo cambia ese movimiento a lo largo del año. A continuación discutimos los resultados principales y sus limitaciones.

\subsection{El error del GPS}

Para un animal que se desplaza unos pocos metros cada 15 minutos, el error del GPS es del mismo orden que el movimiento que se quiere medir. La
caracterización del Capítulo~\ref{Capitulo_1} muestra que este error no es gaussiano ni independiente: la diferencia entre dos posiciones de una tortuga
quieta sigue una t de Student con $\nu \approx 2$, de colas mucho más pesadas que una gaussiana, y el error tiene dos escalas de memoria, una rápida de
unos 30 minutos y una lenta que suma unos $1.3\,\mathrm{m}$ por eje en una a dos horas (Secciones~\ref{I:ruido_noche} y~\ref{I:corr_error}).

Estas características tienen consecuencias directas sobre el modelado. Sin el error, entre el $76$ y el $80 \%$ de los pasos simulados de 15 minutos
son nulos, mientras que en los datos son menos del $0.7 \%$ (Sección~\ref{II:Agregamos_ruido}). Además, la memoria del error hace que una tortuga quieta
parezca alejarse más de $10\,\mathrm{m}$ de su posición inicial: con el error correlacionado y la deriva lenta, cerca del $23 \%$ de los días simulados
de reposo no se clasifica como reposo (Tabla~\ref{tab:Clasif_reposo_err}). Por eso, en lugar de filtrar el error, optamos por incorporarlo al modelo y
aplicar a las simulaciones las mismas métricas y la misma clasificación que a los datos.

\subsection{Un único modelo no alcanza}

El modelo \textit{Stop and Go} con duraciones de los estados en ley de potencias reproduce la forma de ley de potencias del P(d) solo cuando el reposo
decae más lento que la caminata, $\alpha_S < \alpha_W$ (Sección~\ref{II:modelo_unico}), que es lo que se observa en las mediciones del
acelerómetro~\cite{callenbach2025}. Sin embargo, el modelo base no reproduce a la vez el P(d) y el MSD diario. Agregar el error del GPS, la hora de salida
y la variabilidad de la velocidad permite ajustar ambas métricas y la hora de salida con exponentes dentro del rango medido (Secciones~\ref{II:Agregamos_ruido}
a~\ref{II:Var_vel}). Aun así, ese modelo clasifica el $85.4 \%$ de sus días como cambio de refugio, frente al $44.7 \%$ de los datos
(Sección~\ref{II:tipos_de_movs}). Es decir, las métricas promediadas sobre todos los días no alcanzan para restringir el modelo: un modelo puede reproducir
el P(d) y el MSD diario promedio y, al mismo tiempo, generar días muy distintos de los observados.

La clasificación de los días muestra que el movimiento diario es una mezcla de comportamientos con P(d) y MSD diario propios
(Sección~\ref{II:Movimientos_caracterizacion}). Este resultado coincide con lo encontrado en otros animales, donde los modelos con un único modo de
movimiento resultan insuficientes frente a mezclas de varios modos~\cite{morales2004extracting, fryxell2008multiple}. Por eso optamos por un submodelo
para cada tipo de día: solo el error del GPS para el reposo, el \textit{Stop and Go} con todas sus capas para el cambio de refugio y una caminata con
memoria del refugio para la exploración y retorno. Los pesos de la simulación conjunta se eligen para reproducir los porcentajes de cada tipo de día
(Tabla~\ref{tab:Porcentajes_días}), de modo que esos porcentajes se cumplen por construcción. Lo que el modelo conjunto reproduce sin haber sido ajustado
para ello es la combinación del P(d), el MSD diario y la hora de salida de todos los días.

\subsection{Reposo y exploración y retorno}

El peso del submodelo de reposo en la simulación conjunta es de solo el $3.6 \%$, aunque el $22.6 \%$ de los días de los datos se clasifica como reposo.
Esto ocurre porque el $36.1 \%$ de los días del submodelo de exploración y retorno se clasifica como reposo, y el $22.8 \%$ de los días del submodelo
de reposo, como exploración y retorno (Tabla~\ref{tab:Porcentajes_días}). Con un umbral de $10\,\mathrm{m}$, cercano al error del GPS, el reposo y las
exploraciones cortas no se distinguen bien, y lo que en los datos aparece como dos categorías podría ser un mismo comportamiento con excursiones de
distinto alcance. El análisis de sensibilidad del Apéndice~\ref{A:Sensibilidad_10m_clasificacion} apunta en el mismo sentido: al subir el umbral, solo
pasan días de exploración y retorno a reposo. El cambio de refugio, en cambio, aparece como un comportamiento bien definido: su fracción casi no depende
del umbral y su submodelo clasifica el $94.5 \%$ de sus días como cambio de refugio.

\subsection{Hora de salida y velocidad}

La hora de salida del refugio no sigue una distribución exponencial: la tasa de salida crece a lo largo de la mañana y se satura en unos
$0.5\,\mathrm{h^{-1}}$ (Sección~\ref{II:hr_comienzo}). Además, la salida se adelanta en la época de reproducción y se retrasa en la brumación, y en cada
época los días de cambio de refugio empiezan antes que los de exploración y retorno (Capítulo~\ref{Capitulo_3}). Una posible lectura es que la salida
responde a la temperatura y que los días de cambio de refugio empiezan antes porque requieren más tiempo de caminata, pero los datos de GPS solos no
permiten confirmarlo. También hay que tener en cuenta que la hora detectada es posterior a la real, porque la regla de detección necesita que la tortuga
ya se haya alejado del refugio.

La variabilidad de la velocidad fue necesaria para ajustar el P(d) y el MSD diario a la vez. Como los días simulados son independientes, el modelo no
distingue si esa variabilidad ocurre entre días de una misma tortuga o entre tortugas distintas (Sección~\ref{II:Var_vel}). Del mismo modo, el
Apéndice~\ref{A:T_cte} muestra que modelos con duraciones constantes en uno de los estados también reproducen el P(d) y el MSD diario: con posiciones
cada 15 minutos no es posible distinguir entre esas alternativas, y la elección de leyes de potencia se apoya en las mediciones del
acelerómetro~\cite{callenbach2025}. El GPS y el acelerómetro resultan complementarios: el primero fija la escala espacial del movimiento y el segundo,
su estructura temporal.

\subsection{Variación a lo largo del año}

El uso del espacio cambia fuertemente con la época. El radio de giro mensual supera los $150\,\mathrm{m}$ en la reproducción y baja a entre $13$ y
$30\,\mathrm{m}$ entre mayo y agosto, del orden del error del GPS (Figura~\ref{fig:rog_mes}); la tasa de revisita pasa de $0.64$ en la reproducción a
$0.91$ en la brumación (Figura~\ref{fig:tasa_revisita}); y el cambio de refugio, que domina la reproducción ($65 \%$ de los días), pierde peso frente
al reposo durante la brumación (Sección~\ref{III:Caminatas_ev}). Abril y octubre se comportan como meses de transición, posiblemente porque no todos los
individuos entran y salen de la brumación al mismo tiempo.

La selección de refugios sigue una dinámica de exploración y retorno preferencial similar a la descrita para seres humanos~\cite{song2010modeling}: la
probabilidad de volver a un refugio crece de forma aproximadamente lineal con sus visitas previas ($\alpha = 0.97(2)$), la probabilidad de usar uno
nuevo decae como $S^{-\gamma}$ con $\gamma = 0.22(2)$, compatible con el valor de los seres humanos, y el número de refugios distintos crece como
$n^{\mu}$ con $\mu = 0.80(2)$, cercano a la predicción $\mu = 1/(1+\gamma)$ (Sección~\ref{III:refugios}). Con seis individuos y un año de datos, estos
exponentes deben tomarse con cautela.

\subsection{Limitaciones}

Los resultados tienen las limitaciones propias de los datos: seis individuos (cinco en la clasificación diaria) de una sola población, un año de registro,
noches sin mediciones durante parte del año y una brumación profunda medida solo entre las 11:00 y las 15:40 h. El modelo, por su parte, supone un
espacio homogéneo, individuos independientes y días independientes entre sí. La Figura~\ref{fig:Satelite} muestra que el espacio no es homogéneo: durante
la reproducción, las tortugas recorren los bordes de los caminos. Por último, la clasificación depende de un umbral fijo de $10\,\mathrm{m}$, cuya
elección analizamos en el Apéndice~\ref{A:Sensibilidad_10m_clasificacion}.

\section{Conclusiones}\label{IV:conclusiones}

\begin{enumerate}
    \item El error de posición de los GPS utilizados sigue una distribución t de Student con $\nu \approx 2$ y tiene memoria en dos escalas, de unos
    30 minutos y de algunas horas. Para animales tan lentos como \textit{Chelonoidis chilensis}, este error es del orden del movimiento a 15 minutos y
    debe incluirse en el modelo, aplicando a las simulaciones las mismas mediciones que a los datos.
    \item Un modelo \textit{Stop and Go} con duraciones en ley de potencias, en el que el reposo decae más lento que la caminata, con el error del GPS,
    una hora de salida sigmoide y una velocidad lognormal que varía entre días, reproduce el P(d) y el MSD diario de todo el año, pero no la proporción de
    tipos de día observada.
    \item El movimiento diario es una mezcla de tres comportamientos: reposo, exploración y retorno, y cambio de refugio. Un submodelo para cada uno,
    combinados con pesos, reproduce simultáneamente el P(d), el MSD diario, la hora de salida y la clasificación de los días. Con un umbral de
    $10\,\mathrm{m}$, el reposo y las exploraciones cortas no se separan con claridad, mientras que el cambio de refugio sí.
    \item La hora de salida del refugio sigue una función logística (centro 11:19~h, ancho 2~h) con una tasa de salida que crece durante la mañana, y
    cambia con la época y con el tipo de día.
    \item El comportamiento varía fuertemente a lo largo del año: en la reproducción las tortugas usan más espacio, cambian más de refugio y salen más
    temprano que en la brumación. La selección de refugios sigue una dinámica de exploración y retorno preferencial, con $\gamma = 0.22(2)$ y
    $\mu = 0.80(2)$.
\end{enumerate}

\section{Trabajo futuro}\label{IV:futuro}

\begin{itemize}
    \item \textbf{Simulación del año completo.} Combinar los tres submodelos con pesos y horas de salida propias de cada época permitiría simular un año
    y compararlo con las métricas del Capítulo~\ref{Capitulo_3}. Es posible que un único juego de parámetros por submodelo no alcance y que haga falta
    ajustar los parámetros de cada época.
    \item \textbf{Memoria entre días.} En el modelo actual los días son independientes. Incorporar la selección de refugios del Capítulo~\ref{Capitulo_3}
    (retorno preferencial y exploración de refugios nuevos) conectaría el movimiento diario con el uso del espacio en escalas de semanas a meses.
    \item \textbf{Clasificación de los días.} Reemplazar el umbral fijo por una clasificación que tenga en cuenta el modelo del error del GPS, o unificar
    el reposo y las exploraciones cortas en una sola categoría.
    \item \textbf{Factores ambientales.} Incorporar la temperatura horaria permitiría probar si la hora de salida y la entrada en brumación responden a
    un umbral térmico.
    \item \textbf{Heterogeneidad espacial.} Incluir caminos y vegetación a partir de imágenes satelitales, ya que durante la reproducción las tortugas
    usan los bordes de los caminos.
    \item \textbf{Más individuos y dispositivos combinados.} Más individuos permitirían separar la variabilidad de la velocidad entre días y entre
    tortugas, y estudiar diferencias entre sexos. Registrar GPS y acelerómetro en el mismo animal permitiría validar la clasificación y los exponentes
    de las duraciones de los estados.
\end{itemize}

%%% Local Variables:
%%% mode: latex
%%% TeX-master: "template"
%%% End: